\section*{Supplementary Note 7. System-by-system summary}

This note summarises, for each of the 34 benchmark problems, what the model describes, which parameters the method selected and what they mean in the model, how well the reduced model reproduces the full one, and how it predicted held-out data. The meaning of each selected parameter was taken from the reactions and rules of the SBML model in which it appears (following assignment rules and condition-specific parameters to the reactions they affect). The numbers are those of Table~1 of the main text and Supplementary Table~S1: $p$ is the number of estimated parameters, $\kappa$ and VIF are the condition number and largest variance inflation factor of the selected sensitivity columns, $R_{\mathrm{var}}$ is the fraction of sensitivity energy kept, $\cos\Delta$ and $e_{\mathrm{fit}}$ are medians over 15 random scenarios at $\pm$5\% (threshold $\tau_e=0.436$), the level of guarantee is defined in the Results of the main text, and the robustness patterns A--D are those of Fig.~4 of the main text. In the external validation, both the reduced and the full model are refitted to the training data, and the test $\chi^2/n$ (mean squared residual divided by the recorded noise variance; about 1 for a model that predicts within the noise) is computed on the held-out data. The published estimate $\\theta_0$ was obtained with all the data, including the test data, so its test $\\chi^2/n$ is a reference rather than a competing prediction.

\subsection*{Alkan\_SciSignal2018}

Cell-cycle progression, DNA double-strand-break damage and repair under chemotherapy (36 species). Source: \href{https://doi.org/10.1126/scisignal.aat0229}{\nolinkurl{doi:10.1126/scisignal.aat0229}}. The full model has 44 parameters, 10 of which are measurement calibration parameters (scales, offsets or noise levels) that are always re-estimated. The method selected two biological parameters, \texttt{kt} and \texttt{k\_rep\_nhej} (stage~1, ranking by sensitivity energy); the reduced problem has 73\% fewer free parameters than the full one. \texttt{kt} is the cell-cycle transition rate, shared by the transitions G1$\to$S, S$\to$G2 and G2$\to$division; \texttt{k\_rep\_nhej} is the rate of repair by non-homologous end joining, which returns damaged S and G2 cells to the cycle. For this subset $\kappa=1.3$ and VIF $1.1$ (no collinearity), $R_{\mathrm{var}}=94\%$, median $\cos\Delta=0.98$ and median $e_{\mathrm{fit}}=12.1\%$ at $\pm$5\%: the reduced model reproduces the full one, with guarantee level~3 ($\cos\Delta\ge0.9$ and $R_{\mathrm{var}}\ge0.89$); robustness pattern~A: admissible at all perturbation sizes up to $\pm$50\%.

External validation on held-out experimental conditions (555 test and 1178 training measurements): test $\chi^2/n$ 2.75 for the reduced model, 3.38 for the full model refitted to the same data and 2.47 for the published $\theta_0$. The reduced model generalises better. The model was not verified in Lean because the automatic checks of the domain and quasi-positivity conditions fail.

\subsection*{Armistead\_CellDeathDis2024}

Ceramide and sphingolipid metabolism (4 species). Source: \href{https://doi.org/10.1038/s41419-024-07134-2}{\nolinkurl{doi:10.1038/s41419-024-07134-2}}. The full model has 14 parameters. The method selected two biological parameters, \texttt{k\_d} and \texttt{alpha\_cer} (stage~2, ranking by finite-perturbation effects); backward elimination had removed \texttt{k00} and \texttt{k2} from the initial selection; the reduced problem has 86\% fewer free parameters than the full one. \texttt{k\_d} is the degradation rate of ceramide; \texttt{alpha\_cer} is the factor of the ceramide feedback on the de novo production of sphinganine. For this subset $\kappa=1.0$ and VIF $1.0$ (no collinearity), $R_{\mathrm{var}}=40\%$, median $\cos\Delta=1.00$ and median $e_{\mathrm{fit}}=11.7\%$ at $\pm$5\%: the reduced model reproduces the full one, with guarantee level~2 ($\cos\Delta\ge0.9$ as well); robustness pattern~B: admissible up to at least $\pm$20\% but not for larger changes.

External validation on held-out experimental conditions (29 test and 29 training measurements): test $\chi^2/n$ 1.22 for the reduced model, 1.37 for the full model refitted to the same data and 1.16 for the published $\theta_0$. The reduced model generalises better. The hypotheses on the dynamics were verified in Lean (Supplementary Note~4).

\subsection*{Bachmann\_MSB2011}

JAK2/STAT5 signalling induced by erythropoietin (Epo) with negative feedback by CIS and SOCS3 (25 species). Source: \href{https://doi.org/10.1038/msb.2011.50}{\nolinkurl{doi:10.1038/msb.2011.50}}. The full model has 113 parameters, 74 of which are measurement calibration parameters (scales, offsets or noise levels) that are always re-estimated. The method selected two biological parameters, \texttt{CISRNADelay} and \texttt{JAK2ActEpo} (stage~2, ranking by finite-perturbation effects); the reduced problem has 33\% fewer free parameters than the full one. \texttt{JAK2ActEpo} is the rate of Epo-induced activation of JAK2 at the receptor (EpoR:JAK2 $\to$ EpoR:pJAK2); \texttt{CISRNADelay} is the transit rate through the chain of delay compartments of CIS mRNA, i.e.\ it sets the delay of the CIS feedback. For this subset $\kappa=2.8$ and VIF $2.5$ (no collinearity), $R_{\mathrm{var}}=36\%$, median $\cos\Delta=0.85$ and median $e_{\mathrm{fit}}=33.7\%$ at $\pm$5\%: the reduced model reproduces the full one, with guarantee level~1 ($e_{\mathrm{fit}}\le\tau_e$ only); robustness pattern~A: admissible at all perturbation sizes up to $\pm$50\%.

External validation on held-out experimental conditions (202 test and 339 training measurements): test $\chi^2/n$ 4.77 for the reduced model, 12.23 for the full model refitted to the same data and 1.05 for the published $\theta_0$. The reduced model generalises better. The hypotheses on the dynamics were verified in Lean (Supplementary Note~4).

\subsection*{Beer\_MolBioSystems2014}

Growth of bacterial cultures and production of the pigment indigoidine by engineered synthetase variants (4 species, 19 conditions). Source: \href{https://doi.org/10.1039/c3mb70594c}{\nolinkurl{doi:10.1039/c3mb70594c}}. The full model has 72 parameters. \textbf{Not admissible.} The selection stopped at 18 parameters (synthesis, dimerisation and degradation rates of particular strains, growth rates $\beta$ of particular experiments and the initial amount of bacteria), because any further parameter violated the conditioning test ($\kappa=7.2$, VIF $9.998$, just below the limit of 10). Re-estimating these 18 parameters gave a median $e_{\mathrm{fit}}$ of 50\%, above the threshold of 43.6\%; the outputs depend on many strain- and experiment-specific parameters, and no small subset can compensate for all of them. The Lean verification of this model is only conditional, because the growth check fails.

\subsection*{Bertozzi\_PNAS2020}

SIR model of the COVID-19 epidemic, fitted to two regions (3 species). Source: \href{https://doi.org/10.1073/pnas.2006520117}{\nolinkurl{doi:10.1073/pnas.2006520117}}. The full model has 8 parameters. The method selected two biological parameters, \texttt{R0\_NY} and \texttt{I0\_NY} (stage~2, ranking by finite-perturbation effects); the reduced problem has 75\% fewer free parameters than the full one. \texttt{R0\_NY} is the basic reproduction number in New York, which sets the infection rate; \texttt{I0\_NY} is the initial number of infected individuals in New York. For this subset $\kappa=3.9$ and VIF $4.3$ (no collinearity), $R_{\mathrm{var}}=78\%$, median $\cos\Delta=0.92$ and median $e_{\mathrm{fit}}=32.9\%$ at $\pm$5\%: the reduced model reproduces the full one, with guarantee level~2 ($\cos\Delta\ge0.9$ as well); robustness pattern~B: admissible up to at least $\pm$20\% but not for larger changes.

External validation was not possible: both selected parameters belong only to the held-out condition, so they cannot be estimated from the training data. The hypotheses on the dynamics were verified in Lean (Supplementary Note~4).

\subsection*{Blasi\_CellSystems2016}

Acetylation of the four lysines of histone H4 (16 acetylation states). Source: \href{https://doi.org/10.1016/j.cels.2016.01.002}{\nolinkurl{doi:10.1016/j.cels.2016.01.002}}. The full model has 9 parameters. The method selected two biological parameters, \texttt{a\_basal} and \texttt{a\_k16\_k12k16} (stage~1, ranking by sensitivity energy); backward elimination had removed \texttt{a\_k12k16\_k8k12k16}, \texttt{a\_k5k12\_k5k8k12} and \texttt{a\_k8k12k16\_4ac} from the initial selection; the reduced problem has 78\% fewer free parameters than the full one. \texttt{a\_basal} is the basal acetylation rate, from which most of the transition rates are derived by assignment rules; \texttt{a\_k16\_k12k16} is the rate of acetylation of K12 in H4 already acetylated at K16. For this subset $\kappa=1.8$ and VIF $1.4$ (no collinearity), $R_{\mathrm{var}}=79\%$, median $\cos\Delta=0.91$ and median $e_{\mathrm{fit}}=40.2\%$ at $\pm$5\%: the reduced model reproduces the full one, although only just within the threshold, with guarantee level~2 ($\cos\Delta\ge0.9$ as well); robustness pattern~B: admissible up to at least $\pm$20\% but not for larger changes.

External validation was not possible: the problem has a single condition and a single measurement time, so the data cannot be split. The hypotheses on the dynamics were verified in Lean (Supplementary Note~4).

\subsection*{Boehm\_JProteomeRes2014}

Phosphorylation, dimerisation and nuclear import of STAT5A and STAT5B (8 species). Source: \href{https://doi.org/10.1021/pr5006923}{\nolinkurl{doi:10.1021/pr5006923}}. The full model has 9 parameters. The method selected two biological parameters, \texttt{k\_imp\_hetero} and \texttt{k\_phos} (stage~2, ranking by finite-perturbation effects); the reduced problem has 78\% fewer free parameters than the full one. \texttt{k\_phos} is the rate of phosphorylation and dimerisation of STAT5 (the three reactions forming pApA, pApB and pBpB); \texttt{k\_imp\_hetero} is the rate of nuclear import of the heterodimer pApB. For this subset $\kappa=1.0$ and VIF $1.0$ (no collinearity), $R_{\mathrm{var}}=68\%$, median $\cos\Delta=0.93$ and median $e_{\mathrm{fit}}=9.4\%$ at $\pm$5\%: the reduced model reproduces the full one, with guarantee level~2 ($\cos\Delta\ge0.9$ as well); robustness pattern~A: admissible at all perturbation sizes up to $\pm$50\%.

External validation on later time points ($t\ge80$; 18 test and 30 training measurements): test $\chi^2/n$ 2.57 for the reduced model, 11.35 for the full model refitted to the same data and 0.34 for the published $\theta_0$. The reduced model generalises better; the full model fits the training data closely ($\chi^2/n$ 1.34) but predicts the test data poorly, a sign of overfitting. The hypotheses on the dynamics were verified in Lean (Supplementary Note~4).

\subsection*{Borghans\_BiophysChem1997}

Intracellular Ca$^{2+}$ oscillations by Ca$^{2+}$-induced Ca$^{2+}$ release (3 species). Source: \href{https://doi.org/10.1016/s0301-4622(97)00010-0}{\nolinkurl{doi:10.1016/s0301-4622(97)00010-0}}. The full model has 23 parameters, 2 of which are measurement calibration parameters (scales, offsets or noise levels) that are always re-estimated. The method selected two biological parameters, \texttt{Kz} and \texttt{K\_par} (stage~2, ranking by finite-perturbation effects); the reduced problem has 83\% fewer free parameters than the full one. \texttt{Kz} is the threshold constant of the release of Ca$^{2+}$ from the Ca$^{2+}$-sensitive pool into the cytosol; \texttt{K\_par} is the constant of the removal of cytosolic Ca$^{2+}$. For this subset $\kappa=2.7$ and VIF $2.4$ (no collinearity), $R_{\mathrm{var}}=21\%$, median $\cos\Delta=0.86$ and median $e_{\mathrm{fit}}=4.3\%$ at $\pm$5\%: the reduced model reproduces the full one, with guarantee level~1 ($e_{\mathrm{fit}}\le\tau_e$ only); robustness pattern~A: admissible at all perturbation sizes up to $\pm$50\%.

External validation on later time points ($t\ge6.48875$; 37 test and 74 training measurements): test $\chi^2/n$ 0.74 for the reduced model, 0.74 for the full model refitted to the same data and 0.74 for the published $\theta_0$. The two refitted models predict equally well. The hypotheses on the dynamics were verified in Lean (Supplementary Note~4).

\subsection*{Brannmark\_JBC2010}

Insulin signalling in adipocytes: insulin receptor and IRS1 phosphorylation and internalisation (9 species). Source: \href{https://doi.org/10.1074/jbc.M110.106849}{\nolinkurl{doi:10.1074/jbc.M110.106849}}. The full model has 22 parameters, 3 of which are measurement calibration parameters (scales, offsets or noise levels) that are always re-estimated. The method selected two biological parameters, \texttt{k1f} and \texttt{km2} (stage~2, ranking by finite-perturbation effects); the reduced problem has 77\% fewer free parameters than the full one. \texttt{k1f} is the rate of dephosphorylation and recycling of the internalised phosphorylated receptor (IRiP $\to$ IRi); \texttt{km2} is the rate of dephosphorylation of IRS1 (IRSiP $\to$ IRS). For this subset $\kappa=1.9$ and VIF $1.4$ (no collinearity), $R_{\mathrm{var}}=32\%$, median $\cos\Delta=0.85$ and median $e_{\mathrm{fit}}=7.8\%$ at $\pm$5\%: the reduced model reproduces the full one, with guarantee level~1 ($e_{\mathrm{fit}}\le\tau_e$ only); robustness pattern~A: admissible at all perturbation sizes up to $\pm$50\%.

External validation on held-out experimental conditions (26 test and 17 training measurements): test $\chi^2/n$ 5.79 for the reduced model, 2,145 for the full model refitted to the same data and 0.94 for the published $\theta_0$. The reduced model generalises better; the full model fits the training data closely ($\chi^2/n$ 0.36) but predicts the test data poorly, a sign of overfitting. The hypotheses on the dynamics were verified in Lean (Supplementary Note~4).

\subsection*{Bruno\_JExpBot2016}

Enzymatic cleavage of carotenoids ($\beta$-carotene, $\beta$-cryptoxanthin, zeaxanthin) into apocarotenoids (7 species, 6 experiments). Source: \href{https://doi.org/10.1093/jxb/erw356}{\nolinkurl{doi:10.1093/jxb/erw356}}. The full model has 13 parameters. The method selected five biological parameters, \texttt{kb1}, \texttt{init\_bcar2}, \texttt{init\_zea\_1}, \texttt{init\_bcar1} and \texttt{init\_bcry\_1} (stage~1, ranking by sensitivity energy); backward elimination had removed \texttt{szea} and \texttt{init\_b10\_1} from the initial selection; the reduced problem has 62\% fewer free parameters than the full one. \texttt{kb1} is the cleavage rate of $\beta$-carotene; the other four are initial amounts of the substrates in particular experiments: $\beta$-carotene (\texttt{init\_bcar1}, \texttt{init\_bcar2}), zeaxanthin (\texttt{init\_zea\_1}) and $\beta$-cryptoxanthin (\texttt{init\_bcry\_1}). For this subset $\kappa=2.3$ and VIF $1.6$ (no collinearity), $R_{\mathrm{var}}=75\%$, median $\cos\Delta=0.90$ and median $e_{\mathrm{fit}}=37.8\%$ at $\pm$5\%: the reduced model reproduces the full one, with guarantee level~2 ($\cos\Delta\ge0.9$ as well); robustness pattern~A: admissible at all perturbation sizes up to $\pm$50\%.

External validation on held-out experimental conditions (28 test and 49 training measurements): test $\chi^2/n$ 0.91 for the reduced model, 7.86 for the full model refitted to the same data and 0.91 for the published $\theta_0$. The reduced model generalises better; the full model fits the training data closely ($\chi^2/n$ 0.81) but predicts the test data poorly, a sign of overfitting. The hypotheses on the dynamics were verified in Lean (Supplementary Note~4).

\subsection*{Chen\_MSB2009}

ErbB receptor network (ErbB1--4) with the MAPK and PI3K/AKT cascades (500 species). Source: \href{https://doi.org/10.1038/msb.2008.74}{\nolinkurl{doi:10.1038/msb.2008.74}}. The full model has 155 parameters, 3 of which are measurement calibration parameters (scales, offsets or noise levels) that are always re-estimated. The method selected two biological parameters, \texttt{k106} and \texttt{k44} (stage~2, ranking by finite-perturbation effects); the reduced problem has 97\% fewer free parameters than the full one. \texttt{k106} is the binding rate of PIP2 to the receptor complexes that carry PI3K; \texttt{k44} is the binding rate of MEK to phosphorylated Raf, the first step of MEK activation. For this subset $\kappa=1.0$ and VIF $1.0$ (no collinearity), $R_{\mathrm{var}}=25\%$, median $\cos\Delta=0.90$ and median $e_{\mathrm{fit}}=8.2\%$ at $\pm$5\%: the reduced model reproduces the full one, with guarantee level~1 ($e_{\mathrm{fit}}\le\tau_e$ only); robustness pattern~A: admissible at all perturbation sizes up to $\pm$50\%.

External validation on held-out experimental conditions (30 test and 90 training measurements): test $\chi^2/n$ 20.77 for the reduced model (published $\theta_0$: 15.22); the refit of the full model is still running. The hypotheses on the dynamics were verified in Lean (Supplementary Note~4).

\subsection*{Crauste\_CellSystems2017}

CD8 T-cell response to infection: naive, early and late effector and memory cells and pathogen (5 species). Source: \href{https://doi.org/10.1016/j.cels.2017.01.014}{\nolinkurl{doi:10.1016/j.cels.2017.01.014}}. The full model has 12 parameters. \textbf{Not admissible.} The published parameters lie at the edge of a regime in which the solution grows without bound in finite time: changing a single parameter by 1\% (for example, decreasing the proliferation rate of early effectors $\rho_E$) makes the effector populations grow from about $5\times10^2$ to $5\times10^5$ cells within four days, after which no solution exists on the measured horizon. This happened for 8--9 of the 24 one-at-a-time perturbations of $\pm$1\% and for 7 of the 15 scenarios at $\pm$5\%, so the sensitivity matrix and the admissibility test are not defined and no local reduction method is applicable. Consistently, the existence of solutions could only be proved in Lean on a short horizon ($T\le1$, Supplementary Note~4).

\subsection*{Elowitz\_Nature2000}

The repressilator, a synthetic oscillator of three mutually repressing genes with a GFP reporter (8 species). Source: \href{https://doi.org/10.1038/35002125}{\nolinkurl{doi:10.1038/35002125}}. The full model has 21 parameters, 2 of which are measurement calibration parameters (scales, offsets or noise levels) that are always re-estimated. The method selected two biological parameters, \texttt{n\_Hill} and \texttt{tau\_mRNA} (stage~1, ranking by sensitivity energy); backward elimination had removed \texttt{tau\_prot}, \texttt{tau\_prot\_GFP} and \texttt{eff} from the initial selection; the reduced problem has 81\% fewer free parameters than the full one. \texttt{n\_Hill} is the Hill coefficient of transcriptional repression, shared by the four transcription reactions; \texttt{tau\_mRNA} is the mRNA lifetime, which enters the degradation of the three mRNAs and the translation rates. For this subset $\kappa=2.9$ and VIF $2.6$ (no collinearity), $R_{\mathrm{var}}=70\%$, median $\cos\Delta=0.89$ and median $e_{\mathrm{fit}}=35.8\%$ at $\pm$5\%: the reduced model reproduces the full one, with guarantee level~1 ($e_{\mathrm{fit}}\le\tau_e$ only); robustness pattern~B: admissible up to at least $\pm$20\% but not for larger changes.

External validation on later time points ($t\ge410$; 20 test and 38 training measurements): test $\chi^2/n$ 1.50 for the reduced model, 4.45 for the full model refitted to the same data and 1.38 for the published $\theta_0$. The reduced model generalises better. The hypotheses on the dynamics were verified in Lean (Supplementary Note~4).

\subsection*{Fiedler\_BMCSystBiol2016}

RAF/MEK/ERK cascade with ERK feedback, under different inhibitor treatments (6 species). Source: \href{https://doi.org/10.1186/s12918-016-0319-7}{\nolinkurl{doi:10.1186/s12918-016-0319-7}}. The full model has 22 parameters, 8 of which are measurement calibration parameters (scales, offsets or noise levels) that are always re-estimated. The method selected two biological parameters, \texttt{tau2} and \texttt{k6} (stage~2, ranking by finite-perturbation effects); the reduced problem has 55\% fewer free parameters than the full one. \texttt{tau2} is a time constant of the activation input of RAF; \texttt{k6} is the rate of dephosphorylation of pERK, and it also enters the steady state from which every simulation starts. For this subset $\kappa=1.5$ and VIF $1.2$ (no collinearity), $R_{\mathrm{var}}=42\%$, median $\cos\Delta=0.87$ and median $e_{\mathrm{fit}}=14.7\%$ at $\pm$5\%: the reduced model reproduces the full one, with guarantee level~1 ($e_{\mathrm{fit}}\le\tau_e$ only); robustness pattern~A: admissible at all perturbation sizes up to $\pm$50\%.

External validation on held-out experimental conditions (20 test and 52 training measurements): test $\chi^2/n$ 11.21 for the reduced model, 14.11 for the full model refitted to the same data and 1.22 for the published $\theta_0$. The reduced model generalises better. The hypotheses on the dynamics were verified in Lean (Supplementary Note~4).

\subsection*{Fujita\_SciSignal2010}

EGF-dependent AKT signalling and S6 phosphorylation (9 species). Source: \href{https://doi.org/10.1126/scisignal.2000810}{\nolinkurl{doi:10.1126/scisignal.2000810}}. The full model has 19 parameters, 3 of which are measurement calibration parameters (scales, offsets or noise levels) that are always re-estimated. The method selected two biological parameters, \texttt{init\_AKT} and \texttt{reaction\_6\_k1} (stage~2, ranking by finite-perturbation effects); the reduced problem has 74\% fewer free parameters than the full one. \texttt{init\_AKT} is the initial amount of AKT; \texttt{reaction\_6\_k1} is the rate of release of pAKT and pS6 from the pAKT:S6 complex, i.e.\ of S6 phosphorylation. For this subset $\kappa=2.3$ and VIF $1.9$ (no collinearity), $R_{\mathrm{var}}=36\%$, median $\cos\Delta=0.77$ and median $e_{\mathrm{fit}}=29.4\%$ at $\pm$5\%: the reduced model reproduces the full one, with guarantee level~1 ($e_{\mathrm{fit}}\le\tau_e$ only); robustness pattern~A: admissible at all perturbation sizes up to $\pm$50\%.

External validation on held-out experimental conditions (48 test and 96 training measurements): test $\chi^2/n$ 5.94 for the reduced model, 8.04 for the full model refitted to the same data and 5.63 for the published $\theta_0$. The reduced model generalises better. The hypotheses on the dynamics were verified in Lean (Supplementary Note~4).

\subsection*{Giordano\_Nature2020}

SIDARTHE model of the COVID-19 epidemic in Italy, with parameters that change at the dates of the containment measures (8 species). Source: \href{https://pubmed.ncbi.nlm.nih.gov/32322102/}{PubMed 32322102}. The full model has 50 parameters. The method selected two biological parameters, \texttt{alpha\_4} and \texttt{alpha\_28} (stage~2, ranking by finite-perturbation effects); the reduced problem has 96\% fewer free parameters than the full one. \texttt{alpha\_4} and \texttt{alpha\_28} are the transmission coefficient $\alpha$ (infection by undetected asymptomatic individuals) after the changes of day 4 and day 28. For this subset $\kappa=3.7$ and VIF $3.9$ (no collinearity), $R_{\mathrm{var}}=74\%$, median $\cos\Delta=0.97$ and median $e_{\mathrm{fit}}=5.8\%$ at $\pm$5\%: the reduced model reproduces the full one, with guarantee level~2 ($\cos\Delta\ge0.9$ as well); robustness pattern~A: admissible at all perturbation sizes up to $\pm$50\%.

External validation on later time points ($t\ge30$; 112 test and 201 training measurements): test $\chi^2/n$ $<$0.01 for the reduced model, 0.06 for the full model refitted to the same data and $<$0.01 for the published $\theta_0$. The reduced model generalises better; the full model fits the training data closely ($\chi^2/n$ $<$0.01) but predicts the test data poorly, a sign of overfitting. The hypotheses on the dynamics were verified in Lean (Supplementary Note~4).

\subsection*{Isensee\_JCB2018}

Activation of protein kinase A type II by cAMP in sensory neurons (25 species). Source: \href{https://doi.org/10.1083/jcb.201708053}{\nolinkurl{doi:10.1083/jcb.201708053}}. The full model has 46 parameters, 11 of which are measurement calibration parameters (scales, offsets or noise levels) that are always re-estimated. The method selected two biological parameters, \texttt{kf\_RIIp\_C\_2\_\_RII\_C\_2} and \texttt{xi\_kf\_RII\_C\_2\_\_RII\_2} (stage~1, ranking by sensitivity energy); backward elimination had removed \texttt{ks\_AC\_cAMP}, \texttt{KD\_cAMP}, \texttt{xi\_AC\_cAMP\_Fsk}, \texttt{kf\_RIIp\_cAMP\_C\_2\_\_RIIp\_2} and \texttt{kf\_RII\_2\_\_RII\_C\_2} from the initial selection; the reduced problem has 72\% fewer free parameters than the full one. \texttt{kf\_RIIp\_C\_2\_\_RII\_C\_2} is the rate of dephosphorylation of the regulatory subunit within the holoenzyme (pRII$_2$:C$_2$ $\to$ RII$_2$:C$_2$); \texttt{xi\_kf\_RII\_C\_2\_\_RII\_2} is the factor that scales the dissociation of the catalytic subunits from the phosphorylated holoenzyme. For this subset $\kappa=1.0$ and VIF $1.0$ (no collinearity), $R_{\mathrm{var}}=68\%$, median $\cos\Delta=0.87$ and median $e_{\mathrm{fit}}=7.6\%$ at $\pm$5\%: the reduced model reproduces the full one, with guarantee level~1 ($e_{\mathrm{fit}}\le\tau_e$ only); robustness pattern~A: admissible at all perturbation sizes up to $\pm$50\%.

External validation on held-out experimental conditions (219 test and 468 training measurements): test $\chi^2/n$ 1.08 for the reduced model, 1.08 for the full model refitted to the same data and 1.04 for the published $\theta_0$. The two refitted models predict equally well. The hypotheses on the dynamics were verified in Lean (Supplementary Note~4).

\subsection*{Lang\_PLOSComputBiol2024}

Mammalian cell cycle (E2F, RB, cyclins and CDK inhibitors) in the cytoplasm and nucleus (several hundred species). Source: \href{https://doi.org/10.1371/journal.pcbi.1011151}{\nolinkurl{doi:10.1371/journal.pcbi.1011151}}. The full model has 294 parameters, 16 of which are measurement calibration parameters (scales, offsets or noise levels) that are always re-estimated. The method selected two biological parameters, \texttt{kDeE2f1} and \texttt{kDeCe} (stage~2, ranking by finite-perturbation effects); the reduced problem has 94\% fewer free parameters than the full one. \texttt{kDeE2f1} is the degradation rate of E2F1, shared by all the complexes that contain it; \texttt{kDeCe} is the degradation rate of cyclin E in its complexes with the CDK inhibitors p21 and p27. For this subset $\kappa=5.5$ and VIF $8.1$ (no collinearity), $R_{\mathrm{var}}=17\%$, median $\cos\Delta=0.35$ and median $e_{\mathrm{fit}}=16.1\%$ at $\pm$5\%: the reduced model reproduces the full one, with guarantee level~1 ($e_{\mathrm{fit}}\le\tau_e$ only); robustness pattern~A: admissible at all perturbation sizes up to $\pm$50\%.

External validation on later time points ($t\ge86391.6$; 3200 test and 6400 training measurements): test $\chi^2/n$ 8.98 for the reduced model (published $\theta_0$: 7.33); the refit of the full model is still running. The hypotheses on the dynamics were verified in Lean (Supplementary Note~4).

\subsection*{Laske\_PLOSComputBiol2019}

Influenza A virus replication: nuclear import of viral ribonucleoproteins and viral RNA synthesis. Source: \href{https://doi.org/10.1371/journal.pcbi.1006944}{\nolinkurl{doi:10.1371/journal.pcbi.1006944}}. The full model has 13 parameters, 2 of which are measurement calibration parameters (scales, offsets or noise levels) that are always re-estimated. The method selected two biological parameters, \texttt{k\_imp} and \texttt{k\_syn\_R\_M} (stage~1, ranking by sensitivity energy); backward elimination had removed \texttt{k\_bind\_M1} from the initial selection; the reduced problem has 69\% fewer free parameters than the full one. \texttt{k\_imp} is the rate of nuclear import of viral ribonucleoproteins (cytoplasm $\to$ nucleus); \texttt{k\_syn\_R\_M} is the synthesis rate of viral mRNA, from which the rates of the eight mRNA segments are derived. For this subset $\kappa=1.4$ and VIF $1.1$ (no collinearity), $R_{\mathrm{var}}=84\%$, median $\cos\Delta=0.99$ and median $e_{\mathrm{fit}}=10.1\%$ at $\pm$5\%: the reduced model reproduces the full one, with guarantee level~2 ($\cos\Delta\ge0.9$ as well); robustness pattern~A: admissible at all perturbation sizes up to $\pm$50\%.

External validation on held-out experimental conditions (9 test and 33 training measurements): test $\chi^2/n$ 0.76 for the reduced model, 1.88 for the full model refitted to the same data and 0.44 for the published $\theta_0$. The reduced model generalises better. The hypotheses on the dynamics were verified in Lean (Supplementary Note~4).

\subsection*{Liu\_IFACPapersOnLine2025}

Virus--cell model: uninfected cells with logistic growth, five stages of infected cells and free virus (7 states, written as rate rules). Source: \href{https://doi.org/10.1016/j.ifacol.2025.07.160}{\nolinkurl{doi:10.1016/j.ifacol.2025.07.160}}. The full model has 9 parameters. The method selected two biological parameters, \texttt{kappa} and \texttt{phi} (stage~2, ranking by finite-perturbation effects); the reduced problem has 78\% fewer free parameters than the full one. \texttt{kappa} is the carrying capacity of the logistic growth of the cells; \texttt{phi} is the rate of transition between successive stages of infection. For this subset $\kappa=1.3$ and VIF $1.1$ (no collinearity), $R_{\mathrm{var}}=81\%$, median $\cos\Delta=0.92$ and median $e_{\mathrm{fit}}=26.1\%$ at $\pm$5\%: the reduced model reproduces the full one, with guarantee level~2 ($\cos\Delta\ge0.9$ as well); robustness pattern~A: admissible at all perturbation sizes up to $\pm$50\%.

External validation on held-out experimental conditions (50 test and 50 training measurements): test $\chi^2/n$ 0.95 for the reduced model, 1.02 for the full model refitted to the same data and 0.94 for the published $\theta_0$. The reduced model generalises better. The model was not verified in Lean because the model has events that depend on the state.

\subsection*{Lucarelli\_CellSystems2018}

TGF-$\beta$/SMAD signalling and the expression of its target genes (several dozen species). Source: \href{https://doi.org/10.1016/j.cels.2017.11.010}{\nolinkurl{doi:10.1016/j.cels.2017.11.010}}. The full model has 84 parameters. The method selected ten biological parameters, \texttt{init\_Rec}, \texttt{S2tot}, \texttt{S3tot}, \texttt{S\_dephosphos}, \texttt{S\_phos}, \texttt{geneF\_turn}, \texttt{k\_234}, \texttt{geneK\_turn}, \texttt{S4tot} and \texttt{geneH\_turn} (stage~2, ranking by finite-perturbation effects); the reduced problem has 88\% fewer free parameters than the full one. \texttt{init\_Rec} is the amount of receptor; \texttt{S2tot}, \texttt{S3tot} and \texttt{S4tot} are the total amounts of SMAD2, SMAD3 and SMAD4; \texttt{S\_phos} and \texttt{S\_dephosphos} are the rates of phosphorylation and dephosphorylation of SMAD2/3; \texttt{k\_234} is the rate of formation of the ppSMAD2:ppSMAD3:SMAD4 complex; \texttt{geneF\_turn}, \texttt{geneH\_turn} and \texttt{geneK\_turn} are the turnover rates of three target genes. For this subset $\kappa=6.5$ and VIF $9.8$ (no collinearity), $R_{\mathrm{var}}=64\%$, median $\cos\Delta=0.81$ and median $e_{\mathrm{fit}}=43.5\%$ at $\pm$5\%: the reduced model reproduces the full one, although only just within the threshold, with guarantee level~1 ($e_{\mathrm{fit}}\le\tau_e$ only); robustness pattern~C: irregular, with an error that stays close to the threshold and crosses it at some sizes.

External validation on held-out experimental conditions (598 test and 1157 training measurements): test $\chi^2/n$ 1.64 for the reduced model, 3.22 for the full model refitted to the same data and 1.02 for the published $\theta_0$. The reduced model generalises better. The hypotheses on the dynamics were verified in Lean (Supplementary Note~4).

\subsection*{Okuonghae\_ChaosSolitonsFractals2020}

Model of the COVID-19 epidemic in Lagos, Nigeria, with detected and undetected infected classes (9 species). Source: \href{https://doi.org/10.1016/j.chaos.2020.110032}{\nolinkurl{doi:10.1016/j.chaos.2020.110032}}. The full model has 16 parameters. The method selected two biological parameters, \texttt{transmission\_rate\_effective} and \texttt{gamma\_i} (stage~2, ranking by finite-perturbation effects); the reduced problem has 88\% fewer free parameters than the full one. \texttt{transmission\_rate\_effective} is the effective transmission rate (force of infection, susceptible $\to$ exposed); \texttt{gamma\_i} is the recovery rate of detected infected individuals. For this subset $\kappa=1.3$ and VIF $1.1$ (no collinearity), $R_{\mathrm{var}}=53\%$, median $\cos\Delta=1.00$ and median $e_{\mathrm{fit}}=4.1\%$ at $\pm$5\%: the reduced model reproduces the full one, with guarantee level~2 ($\cos\Delta\ge0.9$ as well); robustness pattern~A: admissible at all perturbation sizes up to $\pm$50\%.

External validation on later time points ($t\ge32$; 32 test and 60 training measurements): test $\chi^2/n$ 46,705 for the reduced model, 122,419 for the full model refitted to the same data and 28,748 for the published $\theta_0$. The reduced model generalises better. Even $\theta_0$ has $\chi^2/n\gg1$, which indicates that the noise levels recorded in the benchmark do not match the scatter of the data. The hypotheses on the dynamics were verified in Lean (Supplementary Note~4).

\subsection*{Oliveira\_NatCommun2021}

Model of the COVID-19 epidemic in Bahia, Brazil, with hospital and intensive-care compartments. Source: \href{https://doi.org/10.1038/s41467-020-19798-3}{\nolinkurl{doi:10.1038/s41467-020-19798-3}}. The full model has 12 parameters. The method selected two biological parameters, \texttt{beta\_0} and \texttt{h\_hosp\_rate} (stage~2, ranking by finite-perturbation effects); the reduced problem has 83\% fewer free parameters than the full one. \texttt{beta\_0} is the transmission rate (susceptible $\to$ exposed); \texttt{h\_hosp\_rate} is the fraction of symptomatic individuals who are hospitalised, which splits the outflow of symptomatic cases between recovery, hospital and ICU. For this subset $\kappa=1.2$ and VIF $1.0$ (no collinearity), $R_{\mathrm{var}}=47\%$, median $\cos\Delta=0.95$ and median $e_{\mathrm{fit}}=12.9\%$ at $\pm$5\%: the reduced model reproduces the full one, with guarantee level~2 ($\cos\Delta\ge0.9$ as well); robustness pattern~A: admissible at all perturbation sizes up to $\pm$50\%.

External validation on later time points ($t\ge40$; 40 test and 80 training measurements): test $\chi^2/n$ 915 for the reduced model, 81,453 for the full model refitted to the same data and 460,396 for the published $\theta_0$. The reduced model generalises better; the full model fits the training data closely ($\chi^2/n$ 8.75) but predicts the test data poorly, a sign of overfitting. Even $\theta_0$ has $\chi^2/n\gg1$, which indicates that the noise levels recorded in the benchmark do not match the scatter of the data. The model was not verified in Lean because the model has events that depend on the parameters.

\subsection*{Perelson\_Science1996}

HIV-1 dynamics in a patient treated with a protease inhibitor (infected cells, infectious and non-infectious virus). Source: \href{https://doi.org/10.1126/science.271.5255.1582}{\nolinkurl{doi:10.1126/science.271.5255.1582}}. The full model has 3 parameters. The method selected two biological parameters, \texttt{c} and \texttt{delta} (stage~1, ranking by sensitivity energy); the reduced problem has 33\% fewer free parameters than the full one. \texttt{c} is the clearance rate of free virions; \texttt{delta} is the death rate of productively infected cells. For this subset $\kappa=1.9$ and VIF $1.5$ (no collinearity), $R_{\mathrm{var}}=100\%$, median $\cos\Delta=1.00$ and median $e_{\mathrm{fit}}=0.0\%$ at $\pm$5\%: the reduced model reproduces the full one, with guarantee level~3 ($\cos\Delta\ge0.9$ and $R_{\mathrm{var}}\ge0.89$); robustness pattern~A: admissible at all perturbation sizes up to $\pm$50\%.

External validation on later time points ($t\ge2.038$; 6 test and 10 training measurements): test $\chi^2/n$ $<$0.01 for the reduced model, $<$0.01 for the full model refitted to the same data and $<$0.01 for the published $\theta_0$. All fits reproduce the data almost exactly, so the comparison is not informative. The model was not verified in Lean because the quasi-positivity check fails for two terms.

\subsection*{Rahman\_MBS2016}

Transmission model with infected classes of different severity and a treated class. Source: \href{https://doi.org/10.1016/j.mbs.2016.07.009}{\nolinkurl{doi:10.1016/j.mbs.2016.07.009}}. The full model has 9 parameters. The method selected two biological parameters, \texttt{infected\_moderate\_transmission\_rate} and \texttt{infected\_normal\_worsen\_rate} (stage~2, ranking by finite-perturbation effects); the reduced problem has 78\% fewer free parameters than the full one. \texttt{infected\_moderate\_transmission\_rate} is the transmission rate of moderately infected individuals, to which the rates of the other classes are tied by assignment rules (it sets the force of infection); \texttt{infected\_normal\_worsen\_rate} is the rate of progression from normal to moderate infection. For this subset $\kappa=5.9$ and VIF $9.3$ (no collinearity), $R_{\mathrm{var}}=60\%$, median $\cos\Delta=0.97$ and median $e_{\mathrm{fit}}=3.7\%$ at $\pm$5\%: the reduced model reproduces the full one, with guarantee level~2 ($\cos\Delta\ge0.9$ as well); robustness pattern~A: admissible at all perturbation sizes up to $\pm$50\%.

External validation on later time points ($t\ge15$; 8 test and 15 training measurements): test $\chi^2/n$ $<$0.01 for the reduced model, 0.03 for the full model refitted to the same data and $<$0.01 for the published $\theta_0$. The reduced model generalises better; the full model fits the training data closely ($\chi^2/n$ $<$0.01) but predicts the test data poorly, a sign of overfitting. The hypotheses on the dynamics were verified in Lean (Supplementary Note~4).

\subsection*{Raia\_CancerResearch2011}

IL-13 signalling through JAK/STAT6 with SOCS3 feedback in Hodgkin lymphoma cells (several dozen parameters, 4 conditions). Source: \href{https://doi.org/10.1158/0008-5472.CAN-10-2987}{\nolinkurl{doi:10.1158/0008-5472.CAN-10-2987}}. The full model has 39 parameters, 5 of which are measurement calibration parameters (scales, offsets or noise levels) that are always re-estimated. The method selected two biological parameters, \texttt{SOCS3mRNA\_production} and \texttt{init\_Rec\_i} (stage~2, ranking by finite-perturbation effects); the reduced problem has 82\% fewer free parameters than the full one. \texttt{SOCS3mRNA\_production} is the production rate of SOCS3 mRNA, which drives the negative feedback; \texttt{init\_Rec\_i} is the initial amount of internalised receptor, which is recycled to the membrane. For this subset $\kappa=5.5$ and VIF $8.1$ (no collinearity), $R_{\mathrm{var}}=69\%$, median $\cos\Delta=1.00$ and median $e_{\mathrm{fit}}=3.3\%$ at $\pm$5\%: the reduced model reproduces the full one, with guarantee level~2 ($\cos\Delta\ge0.9$ as well); robustness pattern~A: admissible at all perturbation sizes up to $\pm$50\%.

External validation on held-out experimental conditions (72 test and 133 training measurements): test $\chi^2/n$ 8.23 for the reduced model, 5.10 for the full model refitted to the same data and 0.42 for the published $\theta_0$. Here the full model predicts better; this is the only problem in which it does. The hypotheses on the dynamics were verified in Lean (Supplementary Note~4).

\subsection*{Raimundez\_PCB2020}

EGFR and MET signalling in two gastric cancer cell lines (MKN1 and HS746T) treated with EGF, HGF and cetuximab. Source: \href{https://doi.org/10.1371/journal.pcbi.1007147}{\nolinkurl{doi:10.1371/journal.pcbi.1007147}}. The full model has 136 parameters, 79 of which are measurement calibration parameters (scales, offsets or noise levels) that are always re-estimated. The method selected two biological parameters, \texttt{kdeg\_membran\_\_MKN1} and \texttt{d\_kimp\_pEGFR\_EGF\_2\_\_MKN1\_2\_HS746T} (stage~2, ranking by finite-perturbation effects); backward elimination had removed \texttt{ka\_MAPK\_\_MKN1} from the initial selection; the reduced problem has 40\% fewer free parameters than the full one. \texttt{kdeg\_membran\_\_MKN1} is the degradation rate of membrane receptors (EGFR, EGFR:EGF, EGFR:cetuximab) in MKN1 cells; \texttt{d\_kimp\_pEGFR\_EGF\_2\_\_MKN1\_2\_HS746T} is the difference between the two cell lines in the internalisation rate of phosphorylated receptor dimers (pEGFR:EGF and pMET). For this subset $\kappa=1.1$ and VIF $1.0$ (no collinearity), $R_{\mathrm{var}}=39\%$, median $\cos\Delta=-0.30$ and median $e_{\mathrm{fit}}=37.2\%$ at $\pm$5\%: the reduced model reproduces the full one, with guarantee level~1 ($e_{\mathrm{fit}}\le\tau_e$ only); robustness pattern~A: admissible at all perturbation sizes up to $\pm$50\%.

External validation on held-out experimental conditions (208 test and 419 training measurements): test $\chi^2/n$ 1.19 for the reduced model, 1.46 for the full model refitted to the same data and 0.98 for the published $\theta_0$. The reduced model generalises better. The hypotheses on the dynamics were verified in Lean (Supplementary Note~4).

\subsection*{SalazarCavazos\_MBoC2020}

Multisite phosphorylation of the epidermal growth factor receptor (EGFR). Source: \href{https://doi.org/10.1091/mbc.E19-09-0548}{\nolinkurl{doi:10.1091/mbc.E19-09-0548}}. The full model has 6 parameters. The method selected two biological parameters, \texttt{ratio\_kpkd\_Y1068\_\_FREE} and \texttt{ratio\_kpkd\_YN\_\_FREE} (stage~2, ranking by finite-perturbation effects); the reduced problem has 67\% fewer free parameters than the full one. \texttt{ratio\_kpkd\_Y1068\_\_FREE} and \texttt{ratio\_kpkd\_YN\_\_FREE} are the ratios of phosphorylation to dephosphorylation rate of tyrosine Y1068 and of the other tyrosines in receptors that are not bound to adaptor proteins. For this subset $\kappa=5.4$ and VIF $7.9$ (no collinearity), $R_{\mathrm{var}}=84\%$, median $\cos\Delta=0.97$ and median $e_{\mathrm{fit}}=18.7\%$ at $\pm$5\%: the reduced model reproduces the full one, with guarantee level~2 ($\cos\Delta\ge0.9$ as well); robustness pattern~A: admissible at all perturbation sizes up to $\pm$50\%.

External validation on held-out experimental conditions (3 test and 15 training measurements): test $\chi^2/n$ 102 for the reduced model, 124 for the full model refitted to the same data and 96.93 for the published $\theta_0$. The reduced model generalises better. Even $\theta_0$ has $\chi^2/n\gg1$, which indicates that the noise levels recorded in the benchmark do not match the scatter of the data. The hypotheses on the dynamics were verified in Lean (Supplementary Note~4).

\subsection*{Schwen\_PONE2014}

Binding and uptake of insulin by hepatocytes through receptors of low and high affinity. Source: \href{https://doi.org/10.1371/journal.pone.0133653}{\nolinkurl{doi:10.1371/journal.pone.0133653}}. The full model has 30 parameters, 15 of which are measurement calibration parameters (scales, offsets or noise levels) that are always re-estimated. The method selected two biological parameters, \texttt{ini\_R1} and \texttt{ini\_R2fold} (stage~2, ranking by finite-perturbation effects); the reduced problem has 43\% fewer free parameters than the full one. \texttt{ini\_R1} is the initial amount of low-affinity receptors, and \texttt{ini\_R2fold} the factor that gives the amount of high-affinity receptors from it; both set the binding and uptake of insulin. For this subset $\kappa=3.1$ and VIF $3.0$ (no collinearity), $R_{\mathrm{var}}=46\%$, median $\cos\Delta=0.93$ and median $e_{\mathrm{fit}}=17.1\%$ at $\pm$5\%: the reduced model reproduces the full one, with guarantee level~2 ($\cos\Delta\ge0.9$ as well); robustness pattern~A: admissible at all perturbation sizes up to $\pm$50\%.

External validation on held-out experimental conditions (19 test and 267 training measurements): test $\chi^2/n$ 8.20 for the reduced model, 8.40 for the full model refitted to the same data and 0.77 for the published $\theta_0$. The reduced model generalises better. The model was not verified in Lean because the quasi-positivity check fails for two terms.

\subsection*{Smith\_BMCSystBiol2013}

Insulin signalling through IRS1, PI3K and serine phosphorylation of IRS1. Source: \href{https://doi.org/10.1186/1752-0509-7-41}{\nolinkurl{doi:10.1186/1752-0509-7-41}}. The full model has 25 parameters, 9 of which are measurement calibration parameters (scales, offsets or noise levels) that are always re-estimated. The method selected two biological parameters, \texttt{k8} and \texttt{kminus7b} (stage~2, ranking by finite-perturbation effects); the reduced problem has 56\% fewer free parameters than the full one. \texttt{k8} is the binding rate of PI3K to tyrosine-phosphorylated IRS1; \texttt{kminus7b} is the rate of removal of the serine phosphorylation of IRS1 (reverse of the inhibitory serine phosphorylation). For this subset $\kappa=3.2$ and VIF $3.1$ (no collinearity), $R_{\mathrm{var}}=41\%$, median $\cos\Delta=0.96$ and median $e_{\mathrm{fit}}=7.8\%$ at $\pm$5\%: the reduced model reproduces the full one, with guarantee level~2 ($\cos\Delta\ge0.9$ as well); robustness pattern~A: admissible at all perturbation sizes up to $\pm$50\%.

External validation on held-out experimental conditions (15 test and 47 training measurements): test $\chi^2/n$ 14,977 for the reduced model, 14,978 for the full model refitted to the same data and 82,921 for the published $\theta_0$. The two refitted models predict equally well. Even $\theta_0$ has $\chi^2/n\gg1$, which indicates that the noise levels recorded in the benchmark do not match the scatter of the data. The model was not verified in Lean because the rate laws contain a non-differentiable maximum.

\subsection*{Sneyd\_PNAS2002}

Gating of the IP$_3$ receptor Ca$^{2+}$ channel (open, shut and inactivated states). Source: \href{https://doi.org/10.1073/pnas.032281999}{\nolinkurl{doi:10.1073/pnas.032281999}}. The full model has 15 parameters. The method selected two biological parameters, \texttt{l4} and \texttt{l2} (stage~2, ranking by finite-perturbation effects); the reduced problem has 87\% fewer free parameters than the full one. \texttt{l2} and \texttt{l4} are rate constants of the Ca$^{2+}$-dependent transitions of the receptor from the rest state to the open and inactivated states. For this subset $\kappa=1.1$ and VIF $1.0$ (no collinearity), $R_{\mathrm{var}}=44\%$, median $\cos\Delta=0.88$ and median $e_{\mathrm{fit}}=38.1\%$ at $\pm$5\%: the reduced model reproduces the full one, with guarantee level~1 ($e_{\mathrm{fit}}\le\tau_e$ only); robustness pattern~A: admissible at all perturbation sizes up to $\pm$50\%.

External validation on held-out experimental conditions (45 test and 90 training measurements): test $\chi^2/n$ 2.22 for the reduced model, 2.55 for the full model refitted to the same data and 1.37 for the published $\theta_0$. The reduced model generalises better. The hypotheses on the dynamics were verified in Lean (Supplementary Note~4).

\subsection*{Weber\_BMC2015}

Feedback between protein kinase D (PKD), the ceramide transfer protein CERT and PI4KIII$\beta$ at the Golgi. Source: \href{https://doi.org/10.1186/s12918-015-0147-1}{\nolinkurl{doi:10.1186/s12918-015-0147-1}}. The full model has 36 parameters, 5 of which are measurement calibration parameters (scales, offsets or noise levels) that are always re-estimated. The method selected two biological parameters, \texttt{s12} and \texttt{a21} (stage~2, ranking by finite-perturbation effects); the reduced problem has 81\% fewer free parameters than the full one. \texttt{s12} is the constant production rate of PKD; \texttt{a21} is the first-order removal rate of PI4KIII$\beta$. For this subset $\kappa=3.7$ and VIF $4.0$ (no collinearity), $R_{\mathrm{var}}=25\%$, median $\cos\Delta=0.41$ and median $e_{\mathrm{fit}}=4.1\%$ at $\pm$5\%: the reduced model reproduces the full one, with guarantee level~1 ($e_{\mathrm{fit}}\le\tau_e$ only); robustness pattern~A: admissible at all perturbation sizes up to $\pm$50\%.

External validation on held-out experimental conditions (59 test and 76 training measurements): test $\chi^2/n$ 1.00 for the reduced model, 226 for the full model refitted to the same data and 0.99 for the published $\theta_0$. The reduced model generalises better; the full model fits the training data closely ($\chi^2/n$ 0.78) but predicts the test data poorly, a sign of overfitting. The hypotheses on the dynamics were verified in Lean (Supplementary Note~4).

\subsection*{Zhao\_QuantBiol2020}

Model of the COVID-19 epidemic in China with unquarantined and quarantined infected individuals, in two stages and three regions. Source: \href{https://pubmed.ncbi.nlm.nih.gov/32219006/}{PubMed 32219006}. The full model has 28 parameters. The method selected three biological parameters, \texttt{R\_Stage\_I\_China}, \texttt{R\_Stage\_II\_Wuhan} and \texttt{gamma\_1\_Stage\_I\_Hubei} (stage~2, ranking by finite-perturbation effects); the reduced problem has 89\% fewer free parameters than the full one. \texttt{R\_Stage\_I\_China} and \texttt{R\_Stage\_II\_Wuhan} are the reproduction numbers of stage~I in China and of stage~II in Wuhan, which set the infection rate; \texttt{gamma\_1\_Stage\_I\_Hubei} is the quarantine rate of infected individuals in Hubei in stage~I. For this subset $\kappa=1.2$ and VIF $1.0$ (no collinearity), $R_{\mathrm{var}}=69\%$, median $\cos\Delta=0.90$ and median $e_{\mathrm{fit}}=29.6\%$ at $\pm$5\%: the reduced model reproduces the full one, with guarantee level~2 ($\cos\Delta\ge0.9$ as well); robustness pattern~A: admissible at all perturbation sizes up to $\pm$50\%.

External validation on held-out experimental conditions (25 test and 57 training measurements): test $\chi^2/n$ 0.35 for the reduced model, 0.35 for the full model refitted to the same data and 0.35 for the published $\theta_0$. The two refitted models predict equally well. The hypotheses on the dynamics were verified in Lean (Supplementary Note~4).

\subsection*{Zheng\_PNAS2012}

Methylation of histone H3 at lysines 27 and 36 (combinatorial methylation states). Source: \href{https://doi.org/10.1073/pnas.1205707109}{\nolinkurl{doi:10.1073/pnas.1205707109}}. The full model has 46 parameters. The method selected two biological parameters, \texttt{k01\_02} and \texttt{k13\_12} (stage~2, ranking by finite-perturbation effects); the reduced problem has 96\% fewer free parameters than the full one. \texttt{k01\_02} is the rate of methylation of K36 from mono- to dimethyl in histones without K27 methylation; \texttt{k13\_12} is the rate of demethylation of K36 from tri- to dimethyl in histones with K27me1. For this subset $\kappa=1.4$ and VIF $1.1$ (no collinearity), $R_{\mathrm{var}}=35\%$, median $\cos\Delta=0.57$ and median $e_{\mathrm{fit}}=30.0\%$ at $\pm$5\%: the reduced model reproduces the full one, with guarantee level~1 ($e_{\mathrm{fit}}\le\tau_e$ only); robustness pattern~A: admissible at all perturbation sizes up to $\pm$50\%.

External validation on later time points ($t\ge10$; 30 test and 30 training measurements): test $\chi^2/n$ 1.58 for the reduced model, 20.60 for the full model refitted to the same data and 1.51 for the published $\theta_0$. The reduced model generalises better; the full model fits the training data closely ($\chi^2/n$ 0.06) but predicts the test data poorly, a sign of overfitting. The hypotheses on the dynamics were verified in Lean (Supplementary Note~4).

